% Generated draft: requires mathematical and source review.
\documentclass[11pt]{article}
\usepackage[margin=25mm]{geometry}
\usepackage{fontspec}
\setmainfont{Latin Modern Roman}
\usepackage{amsmath,amssymb,mathtools,unicode-math}
\setmathfont{Latin Modern Math}
\usepackage{enumitem,needspace,xurl,hyperref}
\setlength{\parindent}{0pt}
\setlength{\parskip}{5pt}
\setlength{\emergencystretch}{3em}
\newcommand{\N}{\mathbb{N}}
\newcommand{\Z}{\mathbb{Z}}
\newcommand{\Q}{\mathbb{Q}}
\newcommand{\R}{\mathbb{R}}
\newcommand{\C}{\mathbb{C}}
\newcommand{\F}{\mathbb{F}}
\newcommand{\A}{\mathbb{A}}
\newcommand{\PP}{\mathbb{P}}
\newcommand{\E}{\mathbb{E}}
\newcommand{\eps}{\varepsilon}
\newcommand{\dd}{\,\mathrm{d}}
\begin{document}
\section*{UnsolvedMath ID: 30004320}
\textbf{Source problem code:} catalogue entry
\subsection{Definitions and mathematical statement}
Let \(k\) be a field and let \(t\) be an indeterminate. The Laurent series field is
\[
k((t))=\left\{\sum_{n\geq N}a_nt^n : N\in\Z,\ a_n\in k\right\},
\]
with \(k\) embedded as the constant series. For a \(k\)-scheme \(X\) and a field extension \(L/k\), write
\[
X(L)=\operatorname{Hom}_k(\operatorname{Spec}L,X)
\]
for its set of \(L\)-rational points.

The source considers an affine algebraic \(k\)-group \(G\) and a homogeneous space \(X\) of \(G\). Geometric homogeneity refers to a \(k\)-action \(G\times_k X\to X\) that becomes transitive over an algebraic closure of \(k\). The source does not specify the precise scheme-theoretic convention for a homogeneous space or for an algebraic group; in particular, it does not state smoothness, reducedness, connectedness, or stabilizer hypotheses. These conventions remain unresolved and are not added as assumptions here.

\textbf{Question.} For every field \(k\), every affine algebraic \(k\)-group \(G\), and every homogeneous \(G\)-space \(X\), with the preceding qualification concerning the intended category, does
\[
X(k((t)))\neq\varnothing
\quad\Longrightarrow\quad
X(k)\neq\varnothing
\]
hold? The supplied statement imposes no restriction on the characteristic of \(k\).

The supplied dated assessment distinguishes known torsor settings from the proposed extension to homogeneous spaces and mentions known characteristic-zero cases. It does not give the exact hypotheses of those results. Those qualifications are retained as supplied context, without asserting a theorem or determining the precise remaining open range.
\subsection{Short English statement}
Must every homogeneous space of an affine algebraic group over a field \(k\) that has a point over \(k((t))\) already have a point over \(k\)? The precise geometric conventions and the scope of the reported known cases require review.
\subsection{Sources}
\begin{itemize}
\item Supplied local registry record: UnsolvedMath permanent ID \(30004320\), title \emph{Descent of Rational Points from Laurent Series Fields}, collection \emph{Oberwolfach Reports Open Problems}, exact source problem code \texttt{catalogue entry}. The record gives proposed year 2019 and names no proposer.
\item Supplied report citation: \emph{Mini-Workshop: Rank One Groups and Exceptional Algebraic Groups} (2020), Oberwolfach Reports, pp.\ 3243--3269, DOI 10.4171/owr/2019/52. \url{https://doi.org/10.4171/owr/2019/52}. The background attributes the problem collection to workshop open-problem discussions documented in Oberwolfach Reports.
\item Florence and Gille, \emph{Residues on affine Grassmannians}. \url{https://doi.org/10.1515/crelle-2021-0007}.
\item Florence, \emph{Residues manuscript}. \url{https://webusers.imj-prg.fr/~mathieu.florence/residus11nov.pdf}.
\item \emph{Oberwolfach Reports, torsors and homogeneous spaces}. \url{https://ems.press/content/serial-article-files/46832}.
\item Literature assessment dated \texttt{2026-08-22}: the supplied curation labels the problem open, reports that the torsor statement is established in several settings, and describes its extension to homogeneous spaces as explicitly open beyond known characteristic-zero cases. This is supplied dated curation derived from the OWR record, not an independently verified assessment. The source itself requests independent verification of time-sensitive status and cited claims.
\item The registry labels the statement \texttt{corrected\_verified}, explaining that mathematical content was preserved while notation, typography, omitted definitions, or duplicate clauses were repaired for self-contained publication. This supplied editorial label is not independent mathematical verification of this draft.
\end{itemize}
\end{document}
